% Volume II: Causal Explosion
% Part IV. Time, Delay, and the Maturation of Evidence
% Chapter 21: The Half-Life of an Initial Narrative
% Spine: proof-spines/ch021-spine.md

\chapter{The Half-Life of an Initial Narrative}
\label{ch:ch021}

Initial stories persist because memory, institutions, and reputation keep feeding them after their first presentation. The chapter models this \textbf{half-life of an initial narrative} by
\[
N(t)=N_0e^{-\delta t}+\int_0^t \rho(s)e^{-\delta(t-s)}\,ds,
\]
\ToyModel\ This reinforcement-decay recursion illustrates narrative persistence without pretending to measure any one real system.
where \(\rho(s)\) represents reinforcement through repetition. Narrative authority therefore decays only partially on its own; it is continually replenished by social and institutional investment.

Correction is not symmetrical with accusation. Later evidence enters an attention environment already structured by the original narrative, so even good counterevidence must fight the installed machinery of paperwork, prestige, memory, and moral convenience. Some of that resistance is reinforced deliberately: an actor benefited by the original story can widen \(\rho(s)\) or otherwise stretch the interval before \(t_s\), which is \textbf{strategic delay} rather than the procedural delay this Part otherwise defends, and it should be named as such whenever it appears. The chapter asks how procedures can shorten the narrative's half-life without destroying coordination, so that revisability remains real rather than merely ceremonial.
